\section{Multiple Timestep Algorithm for Coulomb forces\label{sec:Multiple-Timestep-Algorithm}}

For simulations involving a large number of beads in the calculation
of the long-range Coulomb interactions JETSPIN offers the possibility
of using a multiple timestep algorithm which improves the efficiency.
The method is based on that described by Streett \textit{et al} in
Ref \cite{streett1978multiple} and modified for the special context
of electrospinning process.

In the multiple timestep algorithm a primary cutoff $r_{pcut}$ for
the Coulomb pair interactions is given. The cutoff $r_{pcut}$ is
used to define a standard Verlet neighbor list \cite{allen1989computer}.
Long-range Coulomb forces derived from beads in the neighbor list
are usually larger than those due to remaining (far) beads not contained
in the list. Therefore, it is possible to assess the Coulomb energy
if the forces derived from the beads in the neighbor list are calculated
at every timestep, while those from the remaining beads are calculated
much less frequently, and are merely extrapolated over the interval.
JETSPIN handles this procedure as follows: The Verlet neighbor list
is updated at regular intervals $n_{mult}$. The interval between
updates should be usually of the order of \ensuremath{\sim}50 timesteps
$n_{step}$. Whenever the Verlet neighbor list is updated ($n_{step}=1$),
the force contributions from both the beads of the list and outside
the list are calculated. Then, the Coulomb forces are computed in
total in the subsequent timestep ($n_{step}=2$), Hence, for $n_{step}>2$
the Coulomb force acting on a single bead is assessed as sum of the
forces derived from the beads in the neighbor list calculated explicitly,
and those derived from the outer beads (outside the list) which are
calculated by linear extrapolation of the exact forces obtained in
the first two timesteps of the multi-step interval. As soon as the
condition $n_{step}=n_{mult}$ is achieved, the procedure is restarted
(the neighbor list is again updated). It is readily apparent how this
scheme can lead to a significant saving in execution time.

Note that the accuracy of the algorithm is a function of the multiple
step interval, and decreases as $n_{mult}$ increases. Further, the
neighbor list removes the need to scan over all the beads in the simulation
at every timestep. It is worth to highlight that the algorithm is
not time reversible.

We remark that it is also possible to define in input file a maximum
displacement $\Delta r_{pcut}$ in order to achieve a better control
on the accuracy of the multiple timestep algorithm. The maximum displacement
$\Delta r_{pcut}$ works as follows: whenever a bead has traveled
for a distance larger than $\Delta r_{pcut}$ from its position at
multi-step $n_{mult}=1$, the Verlet neighbor list is updated (even
if $n_{step}<n_{mult}$).
